\section*{Chapter \ref{ch:g12:trees} --- Trees and Binary Search Trees}

\begin{solution}{exo:g12:trees:1}
The \omterm{def:g12:trees:tree}{root} $1$ has children $2$ and $3$; $2$ has children $4$ and $5$; $5$ has $7$ and $8$; $3$ has $6$. \omterm{def:g12:trees:height}{Size} $8$;
\omterm{def:g12:trees:height}{height} $3$, along $1$, $2$, $5$, $7$; four leaves: $4$, $7$, $8$, $6$.
\end{solution}

\begin{solution}{exo:g12:trees:2}
$1, 2, 4, 5, 7, 8, 3, 6$.
\end{solution}

\begin{solution}{exo:g12:trees:3}
\omterm{def:g12:trees:traversal}{Inorder}: $4, 2, 7, 5, 8, 1, 6, 3$. \omterm{def:g12:trees:traversal}{Postorder}: $4, 7, 8, 5, 2, 6, 3, 1$.
\end{solution}

\begin{solution}{exo:g12:trees:4}
$50$ at the \omterm{def:g12:trees:tree}{root}; $30$ on its left and $70$ on its right; $20$ and $40$ below $30$, $60$ and $80$ below $70$; $35$ is
smaller than $50$, larger than $30$, smaller than $40$: the left \omterm{def:g12:trees:tree}{child} of $40$. \omterm{def:g12:trees:height}{Height} $3$, the \omterm{def:g12:trees:height}{depth} of $35$.
\end{solution}

\begin{solution}{exo:g12:trees:5}
\smallskip\noindent\hspace*{2em}\code{def count_at_depth(t, d):}

\noindent\hspace*{4em}\code{if t is None:}

\noindent\hspace*{6em}\code{return 0}

\noindent\hspace*{4em}\code{if d == 0:}

\noindent\hspace*{6em}\code{return 1}

\noindent\hspace*{4em}\code{return count_at_depth(t.left, d - 1) + count_at_depth(t.right, d - 1)}

\smallskip\noindent A \omterm{def:g12:trees:tree}{node} at \omterm{def:g12:trees:height}{depth} $d$ in \code{t} is at \omterm{def:g12:trees:height}{depth} $d - 1$ in one of its \omterm{def:g12:trees:tree}{subtrees}. On the \omterm{def:g12:trees:tree}{tree} of
exercise \ref{exo:g12:trees:4}: $1, 2, 4, 1$ for \omterm{def:g12:trees:height}{depths} $0$ to $3$.
\end{solution}

\begin{solution}{exo:g12:trees:6}
It returns $1$ for a single \omterm{def:g12:trees:tree}{node} and $4$ for the \omterm{def:g12:trees:tree}{tree} of the first figure: it counts the \omterm{def:g12:trees:tree}{nodes} on the longest path,
not the links, so it gives $h + 1$. Both conventions are in use; a text must state its own. For this function
$h' = h + 1$, the bounds read $h' \le n \le 2^{h'} - 1$.
\end{solution}

\begin{solution}{exo:g12:trees:7}
\smallskip\noindent\hspace*{2em}\code{def bst_min(t):}

\noindent\hspace*{4em}\code{while t.left is not None:}

\noindent\hspace*{6em}\code{t = t.left}

\noindent\hspace*{4em}\code{return t.value}

\smallskip\noindent Every key smaller than a \omterm{def:g12:trees:tree}{node}'s is in its left \omterm{def:g12:trees:tree}{subtree}: when there is none, the \omterm{def:g12:trees:tree}{node} holds the
smallest key. The path followed is the leftmost one.
\end{solution}

\begin{solution}{exo:g12:trees:8}
A \omterm{def:g12:trees:tree}{tree} of \omterm{def:g12:trees:height}{height} $h$ holds at most $2^{h+1} - 1$ \omterm{def:g12:trees:tree}{nodes}. $2^6 - 1 = 63 < 100$ but $2^7 - 1 = 127 \ge 100$: the \omterm{def:g12:trees:height}{height} is
at least $6$, and $6$ is reached by filling \omterm{def:g12:trees:height}{depths} $0$ to $5$ ($63$ \omterm{def:g12:trees:tree}{nodes}) and placing $37$ \omterm{def:g12:trees:tree}{nodes} at \omterm{def:g12:trees:height}{depth} $6$.
\end{solution}

\begin{solution}{exo:g12:trees:9}
$50, 30, 70, 20, 40, 60, 80, 35$.
\end{solution}

\begin{solution}{exo:g12:trees:10}
The recursive call returns the \omterm{def:g12:trees:tree}{subtree} with the key added, and the caller throws it away. When the call reaches an
empty \omterm{def:g12:trees:tree}{subtree}, it makes the new \omterm{def:g12:trees:tree}{node} and returns it, but nothing links it to its \omterm{def:g12:trees:tree}{parent}. The repair stores the
result:

\smallskip\noindent\hspace*{2em}\code{if key < t.value:}

\noindent\hspace*{4em}\code{t.left = insert(t.left, key)}

\noindent\hspace*{2em}\code{elif key > t.value:}

\noindent\hspace*{4em}\code{t.right = insert(t.right, key)}

\smallskip\noindent The \code{elif} also leaves a key already present out, as \code{bst.py} does.
\end{solution}

\begin{solution}{exo:g12:trees:11}
\smallskip\noindent\hspace*{2em}\code{def is_bst(t, low=None, high=None):}

\noindent\hspace*{4em}\code{if t is None:}

\noindent\hspace*{6em}\code{return True}

\noindent\hspace*{4em}\code{if low is not None and t.value <= low:}

\noindent\hspace*{6em}\code{return False}

\noindent\hspace*{4em}\code{if high is not None and t.value >= high:}

\noindent\hspace*{6em}\code{return False}

\noindent\hspace*{4em}\code{return is_bst(t.left, low, t.value) and is_bst(t.right, t.value, high)}

\smallskip\noindent Each key must lie strictly between \code{low} and \code{high}, the bounds set by all its
ancestors. Checking only the children accepts \code{Node(5, Node(3, None, Node(6)), Node(8))}: $6$ is larger than its
\omterm{def:g12:trees:tree}{parent} $3$, as a right \omterm{def:g12:trees:tree}{child} should be, but it lies in the left \omterm{def:g12:trees:tree}{subtree} of $5$ and is larger than $5$. A search for
$6$ would go right at the \omterm{def:g12:trees:tree}{root} and miss it.
\end{solution}

\begin{solution}{exo:g12:trees:12}
Let $P(t)$: ``$t$ has \omterm{def:g12:trees:height}{size} $n \le 2^{h+1} - 1$, $h$ its \omterm{def:g12:trees:height}{height}''. For the empty \omterm{def:g12:trees:tree}{tree}, $h = -1$ and
$n = 0 = 2^0 - 1$. For a \omterm{def:g12:trees:tree}{node} with \omterm{def:g12:trees:tree}{subtrees} $L$ and $R$, assume $P(L)$ and $P(R)$. Their \omterm{def:g12:trees:height}{heights} are at most
$h - 1$, so their \omterm{def:g12:trees:height}{sizes} are at most $2^h - 1$ each, and
$n = 1 + n_L + n_R \le 1 + 2(2^h - 1) = 2^{h+1} - 1$. Every \omterm{def:g12:trees:binary-tree}{binary tree} is built from the empty \omterm{def:g12:trees:tree}{tree} by forming
\omterm{def:g12:trees:tree}{nodes}, so $P$ holds for all.
\end{solution}

\begin{solution}{exo:g12:trees:13}
By induction on the \omterm{def:g12:trees:height}{size}. The empty \omterm{def:g12:trees:tree}{tree} gives the empty list, which is sorted. For a \omterm{def:g12:trees:tree}{node} of key $k$, both \omterm{def:g12:trees:tree}{subtrees}
are search \omterm{def:g12:trees:tree}{trees}, since the property holds at each of their \omterm{def:g12:trees:tree}{nodes}; by induction their \omterm{def:g12:trees:traversal}{inorders} are sorted. Every key
of the left list is smaller than $k$, every key of the right list larger. The \omterm{def:g12:trees:traversal}{inorder} is the left list, then $k$, then
the right list: sorted.
\end{solution}

\begin{solution}{exo:g12:trees:14}
\smallskip\noindent\hspace*{2em}\code{def rebuild(pre, ino):}

\noindent\hspace*{4em}\code{if not pre:}

\noindent\hspace*{6em}\code{return None}

\noindent\hspace*{4em}\code{root = pre[0]}

\noindent\hspace*{4em}\code{k = ino.index(root)}

\noindent\hspace*{4em}\code{left = rebuild(pre[1:k + 1], ino[:k])}

\noindent\hspace*{4em}\code{return Node(root, left, rebuild(pre[k + 1:], ino[k + 1:]))}

\smallskip\noindent The \omterm{def:g12:trees:traversal}{preorder} starts with the \omterm{def:g12:trees:tree}{root}. In the \omterm{def:g12:trees:traversal}{inorder}, the \omterm{def:g12:trees:tree}{root} separates the left \omterm{def:g12:trees:tree}{subtree}'s values,
$k$ of them, from the right's. In the \omterm{def:g12:trees:traversal}{preorder}, the \omterm{def:g12:trees:tree}{root} is followed by the $k$ values of the left \omterm{def:g12:trees:tree}{subtree}, then the
right's. Each pair of lists is the \omterm{def:g12:trees:traversal}{preorder} and \omterm{def:g12:trees:traversal}{inorder} of a smaller \omterm{def:g12:trees:tree}{tree}. The values must be distinct, or
\code{index} could find the wrong \omterm{def:g12:trees:tree}{node}.
\end{solution}

\begin{solution}{exo:g12:trees:15}
\omterm{def:g12:trees:height}{Size} $3$: five shapes (a chain going left--left, left--right, right--left or right--right, and the \omterm{def:g12:trees:tree}{root} with two
children). \omterm{def:g12:trees:height}{Size} $4$: fourteen. A \omterm{def:g12:trees:tree}{tree} of \omterm{def:g12:trees:height}{size} $n \ge 1$ is a \omterm{def:g12:trees:tree}{root} with a left \omterm{def:g12:trees:tree}{subtree} of some \omterm{def:g12:trees:height}{size} $k$, from $0$ to
$n - 1$, and a right \omterm{def:g12:trees:tree}{subtree} of \omterm{def:g12:trees:height}{size} $n - 1 - k$; the two choices are independent:

\smallskip\noindent\hspace*{2em}\code{def count_shapes(n):}

\noindent\hspace*{4em}\code{if n == 0:}

\noindent\hspace*{6em}\code{return 1}

\noindent\hspace*{4em}\code{return sum(count_shapes(k) * count_shapes(n - 1 - k) for k in range(n))}

\smallskip\noindent It gives $1, 1, 2, 5, 14, 42$ for $n = 0$ to $5$; for $4$, $5 + 2 + 2 + 5 = 14$.
\end{solution}

\begin{solution}{pb:g12:trees:1}
\textbf{1.} Each comparison moves one level down, along one path from the \omterm{def:g12:trees:tree}{root}; a path in a \omterm{def:g12:trees:tree}{tree} of \omterm{def:g12:trees:height}{height} $h$ has at
most $h + 1$ \omterm{def:g12:trees:tree}{nodes}.

\textbf{2.} Yes: the search for an absent key also follows one path, until it falls off below a \omterm{def:g12:trees:tree}{node}.

\textbf{3.} The \omterm{def:g12:trees:traversal}{inorder} of a search \omterm{def:g12:trees:tree}{tree} is its keys in increasing order (exercise \ref{exo:g12:trees:13}): it is
$1, 2, \dots, 1023$ for every \omterm{def:g12:trees:tree}{tree}, whatever its shape.

\textbf{4.} The first key of the order: a key, once inserted, never moves, and the first one becomes the \omterm{def:g12:trees:tree}{root}.

\textbf{5.} Each new key is larger than all the keys already in the \omterm{def:g12:trees:tree}{tree}: compared with each \omterm{def:g12:trees:tree}{node} on its way, it goes
right every time, and hangs as the right \omterm{def:g12:trees:tree}{child} of the previous key. The \omterm{def:g12:trees:tree}{tree} is one chain of $n$ \omterm{def:g12:trees:tree}{nodes}, of \omterm{def:g12:trees:height}{height}
$n - 1$.

\textbf{6.} For instance $1, 7, 2, 6, 3, 5, 4$: each key is the smallest or the largest of those not yet inserted, so
it hangs below the previous one. ($7, 1, 6, 2, 5, 3, 4$ is the one in the output.)

\textbf{7.} The recursive \code{height} calls itself once per level, $1024$ calls deep including the empty \omterm{def:g12:trees:tree}{tree}, more
than the $1000$ frames allowed. \code{search} is a loop: a deep \omterm{def:g12:trees:tree}{tree} costs it time, not stack.

\textbf{8.} $1023$: the key $1023$ is at the end of the chain, and every \omterm{def:g12:trees:tree}{node} lies on the path to it.

\textbf{9.} $1023 = 2^{10} - 1$ \omterm{def:g12:trees:tree}{nodes} need $2^{h+1} - 1 \ge 1023$, so $h \ge 9$; \omterm{def:g12:trees:height}{height} $9$ is reached by a \omterm{def:g12:trees:tree}{tree} with
every \omterm{def:g12:trees:height}{depth} from $0$ to $9$ full.

\textbf{10.} The middle key first, $512$; then the middles of the two halves, $256$ and $768$; then the middles of the
four quarters, $128$, $384$, $640$, $896$; and so on, level by level. Each key comes after the keys above it in the
full \omterm{def:g12:trees:tree}{tree}, so it lands at its place.

\textbf{11.} $10$, on any path of the full \omterm{def:g12:trees:tree}{tree}.

\textbf{12.} $4, 2, 6, 1, 3, 5, 7$: \omterm{def:g12:trees:height}{height} $2$ in the output, the smallest for $7 = 2^3 - 1$ \omterm{def:g12:trees:tree}{nodes}.

\textbf{13.} Smallest $16$, largest $31$, mean $21.13$.

\textbf{14.} $20$, for $211$ \omterm{def:g12:trees:tree}{trees}. A \omterm{def:g12:trees:height}{height} of $25$ or more: $23 + 15 + 3 + 2 + 3 + 1 = 47$ \omterm{def:g12:trees:tree}{trees} out of $1000$.

\textbf{15.} $21.13 / 9 \approx 2.3$ times the best; $1022 / 21.13 \approx 48$ times smaller than the worst.

\textbf{16.} The chains are the orders where each key is the smallest or the largest of those left: two choices for
each of the first $1022$ keys, so $2^{1022}$ orders, about $10^{308}$. There are $1023!$ orders in all, about
$10^{2637}$. A thousand random orders have no chance worth counting of meeting a chain.

\textbf{17.} $32$, in the \omterm{def:g12:trees:tree}{tree} of \omterm{def:g12:trees:height}{height} $31$; against $1023$ in the sorted \omterm{def:g12:trees:tree}{tree}.

\textbf{18.} So that the run gives the same $1000$ orders each time: the printed output and the chart agree with each
other and with any rerun.

\textbf{19.} The table's order has nothing to do with alphabetical order, so it behaves like a random one: ten titles
gave \omterm{def:g12:trees:height}{height} $5$. Sorted by title, the table would give a chain: \omterm{def:g12:trees:height}{height} $79$ for its $80$ titles, and $80$ \omterm{def:g12:trees:tree}{nodes}
visited for the last.

\textbf{20.} \textbf{For $1023$ keys, the best \omterm{def:g12:trees:height}{height} is $9$ and the worst $1022$; over $1000$ random orders the
measured \omterm{def:g12:trees:height}{heights} ran from $16$ to $31$, with a mean of $21.13$.}
\end{solution}
